Mit $\tau = \tEaO$ sind die mittleren Geschwindigkeiten in der ersten und in der zweiten Sekunde
\begin{align}
 \bar v_1 &= \vEaF = \vEa, \qquad \bar v_2 = \vEbF = \vEb
\end{align}
Bei gleichmässiger Beschleunigung ist die mittlere Geschwindigkeit eines Intervalls gleich der Geschwindigkeit in seiner Mitte, also bei $\tau/2$ bzw. $3\tau/2$. Zwischen diesen Zeitpunkten liegt die Zeit $\tau$, deshalb ist
\begin{align}
 a &= \frac{\bar v_2-\bar v_1}{\tau} = \aEF = \aE
\end{align}
Von der Mitte der ersten Sekunde zurück zum Anfang nimmt die Geschwindigkeit um $a\,\tau/2$ ab:
\begin{align}
 v_0 &= \bar v_1 - a\,\frac{\tau}{2} = \frac{s_1}{\tau}-\frac{s_2-s_1}{2\,\tau} = \vEnF \\
   &= \frac{3\cdot\sEa-\sEb}{2\cdot\tEa} = \vEn \approx \result{\vEnP{0}{2}}
\end{align}
